\section{Exercise 4}

Bước kiểm tra dùng \texttt{sltiu}, tức so sánh không dấu với 16. Số âm ở dạng unsigned
là một số không dấu rất lớn nên bị loại bởi cùng một lệnh với các số $\geq 16$.
Nếu nhập sai thì chương trình quay về nhãn \texttt{input} để nhập lại:

\begin{lstlisting}[language=MIPS,numbers=none]
input:
  li $v0, 4          # syscall 4: in prompt
  la $a0, prompt
  syscall

  li $v0, 5          # syscall 5: đọc số nguyên
  syscall

  sltiu $t0, $v0, 16 # 1 nếu 0 <= n < 16, so sánh không dấu
  beq $t0, 0, input  # ngoài khoảng thì nhập lại
\end{lstlisting}

Phần in dùng \texttt{srlv} để dịch phải theo số bit lấy từ thanh ghi.
Bộ đếm chạy từ 3 xuống 0, tương ứng với 4 bit từ trọng số 8 đến 1.
Mỗi vòng dịch phải, dùng \texttt{andi} lấy bit, cộng \texttt{'0'} rồi in bằng syscall 11.

\begin{lstlisting}[language=MIPS,numbers=none]
  move $s0, $v0        # lưu n vào $s0 trước khi in
  li $t0, 3            # bắt đầu từ bit 3, ứng với trọng số 8

  li $v0, 4            # syscall 4: in chuỗi
  la $a0, output_msg   # "Its binary form is: "
  syscall

print_loop:
  srlv $t1, $s0, $t0   # đưa bit thứ $t0 xuống vị trí 0
  andi $t1, $t1, 1     # lấy đúng bit đó
  addi $t1, $t1, '0'   # đổi sang mã ASCII

  li $v0, 11           # syscall 11: in ký tự
  move $a0, $t1
  syscall

  subi $t0, $t0, 1     # lùi sang bit thấp hơn
  bgez $t0, print_loop # còn bit thì lặp
\end{lstlisting}

Nhập 20 bị từ chối, sau đó nhập 10:

\begin{lstlisting}[language={},numbers=none]
Please enter a positive integer less than 16: 20
Please enter a positive integer less than 16: 10
Its binary form is: 1010
\end{lstlisting}
